\documentclass[11pt]{article}
\usepackage[margin=1in]{geometry}
\usepackage{amsmath,amssymb,amsthm}
\usepackage[colorlinks=true,urlcolor=blue,citecolor=blue]{hyperref}
\newtheorem{theorem}{Theorem}
\newtheorem{lemma}{Lemma}
\newcommand{\C}{\mathbb C}
\newcommand{\N}{\mathbb N}
\newcommand{\Rea}{\operatorname{Re}}
\newcommand{\Ima}{\operatorname{Im}}
\title{A parabolic counterexample to a conjecture on chaotic polynomial shifts}
\author{Principia Math in Progress}
\date{12 September 2026}
\begin{document}
\maketitle
\begin{abstract}
For every \(1<p<\infty\), the polynomial
\(P(x)_j=x_{j+1}+x_{j+1}^2\) on complex \(\ell^p\) is topologically
mixing and has dense periodic points. Its scalar polynomial has derivative
one at zero. This disproves the proposed necessity of a repelling fixed
point in MathDB 394128. The proof uses an explicit inverse branch and an
elementary five-step covering argument.
\end{abstract}

\paragraph{Task 1. Resolve the proposed characterization.}
The conjecture on printed page 2241 of \cite{owr} asks whether, for a complex
polynomial \(q\) of degree greater than one with \(q(0)=0\), the map
\[
 P:\ell^p\longrightarrow\ell^p,\qquad
 P(x)_j=q(x_{j+1}),
\]
is Devaney chaotic exactly when zero is a repelling fixed point of \(q\).
Here \(1\le p<\infty\), and the source defines Devaney chaos as topological
transitivity together with density of periodic points.

\begin{theorem}\label{thm:main}
Let \(q(z)=z+z^2\). For each \(1<p<\infty\), the associated map \(P\)
on complex \(\ell^p\) is a continuous polynomial, is topologically mixing,
has dense periodic points, and has sensitive dependence on initial
conditions. Nevertheless \(q'(0)=1\), so zero is not repelling.
In particular \(p=2\) gives a counterexample to the conjecture.
\end{theorem}

Throughout, superscripts on \(q\) and on the inverse branch below denote
iteration, and \(H=\{z\in\C:\Rea z>0\}\).

\begin{lemma}\label{lem:cover}
For \(q(z)=z+z^2\), one has \(q^5(H)=\C\).
\end{lemma}
\begin{proof}
First,
\begin{equation}\label{eq:parabola}
q(H)=\{s+it:s+t^2>0\}.
\end{equation}
Indeed write a square root of \(w+1/4\) as \(A+iB\), choosing \(A\ge0\).
The two preimages of \(w\) are \(-1/2\pm(A+iB)\). Moreover,
\[
 \Rea w+(\Ima w)^2=(A^2-1/4)(1+4B^2).
\]
Thus a preimage belongs to \(H\) precisely when \(A>1/2\), proving
\eqref{eq:parabola}.

If \(w\notin q^2(H)\), neither of its preimages belongs to \(q(H)\).
Applying \eqref{eq:parabola} to these preimages gives
\(A+B^2\le1/2\). Consequently \(0\le A\le1/2\) and
\[
 |w+1/4|=A^2+B^2\le A^2-A+1/2\le1/2,
 \qquad |w|\le3/4.
\]
We have proved
\begin{equation}\label{eq:outside}
 \{w:|w|>3/4\}\subset q^2(H).
\end{equation}

Put
\[
 v=\frac{-1+i}{2},\qquad
 c=\frac{-1-\sqrt{-1+2i}}2,
\]
using the principal square root. Then \(q(c)=v\), \(q(v)=-1/2\), and
\[
 \Rea\sqrt{-1+2i}
 =\sqrt{\frac{\sqrt5-1}{2}}>1/2.
\]
Hence \(\Rea c<-3/4\), while \((q^3)'(c)=0\).
For every \(w\in\C\), some root \(u\) of \(q^3(z)-w\) must therefore
satisfy \(|u|>3/4\). To justify this last assertion directly, if all
roots \(r_1,\ldots,r_8\), counted with multiplicity, lay in the closed
disk of radius \(3/4\), the logarithmic derivative at \(c\) would give
\[
 0=\frac{c(q^3)'(c)}{q^3(c)-w}
   =\sum_{\nu=1}^8\frac1{1-r_\nu/c}.
\]
The denominator on the left is nonzero under that assumption.
Every summand on the right has positive real part, since \(|r_\nu/c|<1\);
this is a contradiction. By \eqref{eq:outside}, the resulting root
has the form \(u=q^2(h)\), \(h\in H\), and then \(q^5(h)=w\).
\end{proof}

\begin{lemma}\label{lem:inverse}
The map
\[
 g(z)=\frac{-1+\sqrt{1+4z}}2
     =\frac{2z}{1+\sqrt{1+4z}},\qquad z\in H,
\]
with the principal square root, maps \(H\) into itself and satisfies
\(q\circ g=\operatorname{id}_H\). For \(h\in H\) and \(k\ge1\),
\begin{equation}\label{eq:decay}
 |g^k(h)|\le\frac{(1+|h|)^2}{k}.
\end{equation}
\end{lemma}
\begin{proof}
For \(z\in H\), the real part of \(\sqrt{1+4z}\) is greater than one:
its square is
\((|1+4z|+\Rea(1+4z))/2>1\).
Thus \(g(z)\in H\), and direct substitution proves the inverse identity.
Set \(a_k=g^k(h)\) and \(M=|h|\). Since
\(a_k=a_{k+1}(1+a_{k+1})\), the moduli \(|a_k|\) decrease and are at
most \(M\). Also
\[
 \frac1{a_{k+1}}-\frac1{a_k}=\frac1{1+a_{k+1}},\qquad
 \Rea\frac1{1+a_{k+1}}
 =\frac{1+\Rea a_{k+1}}{|1+a_{k+1}|^2}
 \ge\frac1{(1+M)^2}.
\]
Summing gives
\(\Rea(1/a_k)\ge \Rea(1/h)+k/(1+M)^2>k/(1+M)^2\).
Since \(|1/a_k|\ge\Rea(1/a_k)\), this proves \eqref{eq:decay}.
\end{proof}

\begin{proof}[Proof of Theorem \ref{thm:main}]
The coordinatewise product is a continuous bilinear map on \(\ell^p\),
because
\[
 \|(x_jy_j)_j\|_p\le\|x\|_\infty\|y\|_p
                    \le\|x\|_p\|y\|_p.
\]
Composing its associated quadratic polynomial with the bounded backward
shift, and adding that shift, gives precisely \(P\). Thus \(P\) is a
continuous polynomial of degree two and is defined on all of \(\ell^p\).
Induction yields
\begin{equation}\label{eq:iterate}
 (P^N x)_j=q^N(x_{j+N}).
\end{equation}

For mixing, let \(U,V\) be nonempty open sets and choose finitely
supported \(a\in U\), \(b\in V\). Take \(K\) containing both supports.
For each \(j\) with \(b_j\ne0\), Lemma \ref{lem:cover} supplies \(h_j\in H\)
with \(q^5(h_j)=b_j\); put \(C_j=(1+|h_j|)^2\).
For every integer \(N>K\), \(N\ge6\), set
\[
 x^{(N)}=a+\sum_{\{j:b_j\ne0\}}g^{N-5}(h_j)e_{N+j}.
\]
The appended coordinates are disjoint from the support of \(a\).
The inverse identity and \eqref{eq:iterate} give
\(P^N(x^{(N)})=b\), and
\[
 \|x^{(N)}-a\|_p
 \le \frac{\bigl(\sum_{\{j:b_j\ne0\}}C_j^p\bigr)^{1/p}}{N-5}
 \longrightarrow0.
\]
Thus \(P^N(U)\cap V\ne\varnothing\) for every sufficiently large \(N\).
This proves mixing, in fact for every \(1\le p<\infty\).
If hypercyclicity is used in the definition, it follows as well:
for a countable base \((V_m)\), each
\(\bigcup_{n\ge0}P^{-n}(V_m)\) is open and dense, so the Baire theorem
provides a point whose orbit meets every \(V_m\).

For density of periodic points, let \(a\) be supported in
\(\{1,\ldots,K\}\). For each nonzero \(a_j\), choose \(h_j\in H\) with
\(q^5(h_j)=a_j\), and again let \(C_j=(1+|h_j|)^2\).
Given an integer \(N>K\), \(N\ge10\), define
\[
 y_j=a_j\quad(1\le j\le K),\qquad
 y_{j+kN}=g^{kN-5}(h_j)
       \quad(a_j\ne0,\ k\ge1),
\]
and set every remaining coordinate to zero. Residue classes modulo \(N\)
do not overlap. As \(kN-5\ge kN/2\),
\begin{equation}\label{eq:tail}
 \|y-a\|_p^p
 \le \sum_{\{j:a_j\ne0\}}\sum_{k=1}^\infty
                    \frac{C_j^p}{(kN-5)^p}
 \le \left(\frac2N\right)^p
       \left(\sum_{k=1}^\infty k^{-p}\right)
       \sum_{\{j:a_j\ne0\}}C_j^p.
\end{equation}
For \(p>1\) this is finite and tends to zero as \(N\to\infty\).
The first block satisfies \(q^N(y_{j+N})=q^5(h_j)=a_j\).
In every subsequent nonzero residue class,
\[
 q^N\bigl(g^{(k+1)N-5}(h_j)\bigr)=g^{kN-5}(h_j)\quad(k\ge1).
\]
Zero residue classes remain zero, so \eqref{eq:iterate} proves \(P^N y=y\).
Finitely supported vectors are dense; hence periodic points are dense.

For completeness, sensitivity has the uniform constant one. Given
\(x\in\ell^p\) and \(\varepsilon>0\), choose \(0<b<\varepsilon/4\).
The positive real iterates \(q^n(b)\) tend to infinity: they increase,
and a finite limit would be a positive fixed point of \(q\), which
does not exist. Choose \(j\) so large that \(|x_j|<\varepsilon/4\) and
\(q^{j-1}(b)>2\). Replacing coordinate \(j\) of \(x\) by zero or by \(b\)
gives two points within distance \(\varepsilon\) of \(x\).
After \(j-1\) iterates their first coordinates differ by more than two.
The triangle inequality puts at least one more than distance one from
\(P^{j-1}x\). Finally \(q'(0)=1\), completing the counterexample.
\end{proof}

The summability step \eqref{eq:tail} requires \(p>1\). No assertion about
dense periodic points on \(\ell^1\) is made here. The original equivalence
is already false on \(\ell^2\).

\paragraph{Task 2. Determine the parabolic summability threshold.}
For \(q_m(z)=z+z^{m+1}\), \(m\ge1\), determine exactly which exponents
\(1\le p<\infty\) give dense periodic points for the associated shift.
Relate the threshold to the decay of inverse iterates near zero,
including the borderline exponent.

\paragraph{Task 3. Analyze irrationally indifferent fixed points.}
For polynomials with \(q(0)=0\) and \(q'(0)=e^{2\pi i\alpha}\),
\(\alpha\) irrational, determine which local behaviors at zero allow
Devaney chaos on \(\ell^p\). Treat the Siegel and Cremer cases separately.

\paragraph{Task 4. Investigate stronger recurrence.}
Determine whether the shift in Theorem \ref{thm:main} is frequently
hypercyclic and whether it admits a fully supported invariant Borel
probability measure that is mixing.

\paragraph{Task 5. Extend the construction to weighted spaces.}
Find necessary and sufficient weight conditions for this parabolic
polynomial shift to be well defined, mixing, and Devaney chaotic on
weighted \(\ell^p\) and K\"othe echelon spaces.

\begin{thebibliography}{9}
\bibitem{owr}
F. Mart\'inez-Gim\'enez (joint work with A. Peris),
Hypercyclic and chaotic polynomials on Fr\'echet spaces,
in \emph{Mini-Workshop: Hypercyclicity and Linear Chaos},
Oberwolfach Report 37/2006, pp.~2238--2241.
\href{https://ems.press/journals/owr/articles/1323}{Publisher report}.
\end{thebibliography}
\end{document}
